\documentclass[a4paper]{article}
\usepackage{vntex}
%\usepackage[enetglish,vinam]{babel}
%\usepackage[utf8]{inputenc}
\usepackage{mathptmx}[ptm]
%\usepackage[utf8]{inputenc}
%\usepackage[francais]{babel}
\usepackage{a4wide,amssymb,epsfig,latexsym,array,hhline,fancyhdr}
\usepackage[normalem]{ulem}
%\usepackage{soul}
\usepackage{placeins} 
\usepackage[makeroom]{cancel}
\usepackage{amsmath}
\usepackage{amsthm}
\usepackage{multicol,longtable,amscd}
\usepackage{diagbox}%Make diagonal lines in tables
\usepackage{booktabs}
\usepackage{alltt}
\usepackage[framemethod=tikz]{mdframed}% For highlighting paragraph backgrounds
\usepackage{caption,subcaption}
\usepackage{float}
\usepackage{lastpage}
\usepackage[lined,boxed,commentsnumbered]{algorithm2e}
\usepackage{enumerate}
\usepackage{color}
\usepackage{graphicx}							% Standard graphics package
\usepackage{array}
\usepackage{tabularx, caption}
\usepackage{multirow}
\usepackage{multicol}
\usepackage{rotating}
\usepackage{graphics}
\usepackage{geometry}
\usepackage{setspace}
\usepackage{epsfig}
\usepackage{tikz}

\usetikzlibrary{arrows,snakes,backgrounds,calc}
\usepackage[unicode]{hyperref}
\hypersetup{urlcolor=blue,linkcolor=black,citecolor=black,colorlinks=true} 
\usepackage{listings}

%\usepackage{pstcol} 								% PSTricks with the standard color package

\usepackage[normalem]{ulem}

\newtheorem{theorem}{{\bf Định lý}}
\newtheorem{property}{{\bf Tính chất}}
\newtheorem{proposition}{{\bf Mệnh đề}}
\newtheorem{corollary}[proposition]{{\bf Hệ quả}}
\newtheorem{lemma}[proposition]{{\bf Bổ đề}}
\theoremstyle{definition}
\newtheorem{exer}{Bài toán}
\addtocontents{toc}{\protect\thispagestyle{empty}}
% Remove page number in contents page. 
\addtocontents{lof}{\protect\thispagestyle{empty}}
% Remove page number in list of figure page. 
\addtocontents{lot}{\protect\thispagestyle{empty}}
% Remove page number in list of table page. 
\def\thesislayout{	% A4: 210 × 297
	\geometry{
		a4paper,
		total={160mm,240mm},  % fix over page
		left=30mm,
		top=22mm,
            right=20mm,
            bottom=20mm,
	}
}
\def\thesisheadlayout{	% A4: 210 × 297
	\geometry{
		a4paper,
		total={160mm,240mm},  % fix over page
		left=30mm,
		top=10mm,
	}
}
\thesislayout
\lstset{
language=R,
basicstyle=\footnotesize\sffamily,
commentstyle=\ttfamily\color{black},
numbers=left,
numberstyle=\ttfamily\color{black}\footnotesize,
stepnumber=1,
numbersep=5pt,
backgroundcolor=\color{white},
showspaces=false,
showstringspaces=false,
showtabs=false,
frame=single,
tabsize=2,
captionpos=b,
breaklines=true,
breakatwhitespace=false,
title=\lstname,
escapeinside={},
keywordstyle={},
morekeywords={}
}
%\usepackage{fancyhdr}
\setlength{\headheight}{40pt}
\pagestyle{fancy}
\fancyhead{} % clear all header fields
\renewcommand{\footruleskip}{1mm}

\fancyhead[L]{
 \begin{tabular}{rl}
    \begin{picture}(25,15)(0,0)
    \put(0,-8){
\includegraphics[width=10mm, height=10mm]{Images/hcmut.png}}
    %\put(0,-8){
\epsfig{width=10mm,figure=hcmut.eps}}
   \end{picture}&
	%
\includegraphics[width=8mm, height=8mm]{hcmut.png} & %
	\begin{tabular}{l}
		\textbf{  Trường Đại Học Bách Khoa - Đại học Quốc gia TP.HCM }\\
		\textbf{  Khoa Khoa Học \& Kỹ Thuật Máy Tính}
	\end{tabular} 	
 \end{tabular}
}
\fancyhead[R]{
	\begin{tabular}{l}
		\tiny \bf \\
		\tiny \bf 
	\end{tabular}  }
\fancyfoot{} % clear all footer fields
\fancyfoot[L]{\scriptsize  Báo cáo Bài tập lớn Khai phá dữ liệu - HK2 2024 - 2025}
\fancyfoot[R]{\scriptsize  Trang {\thepage}/\pageref{LastPage}}

\renewcommand{\headrulewidth}{0.3pt}
\renewcommand{\footrulewidth}{0.3pt}

%%%
\setcounter{secnumdepth}{4}
\setcounter{tocdepth}{3}
\makeatletter
\newcounter {subsubsubsection}[subsubsection]
\renewcommand\thesubsubsubsection{\thesubsubsection .\@alph\c@subsubsubsection}
\newcommand\subsubsubsection{\@startsection{subsubsubsection}{4}{\z@}%
                                     {-3.25ex\@plus -1ex \@minus -.2ex}%
                                     {1.5ex \@plus .2ex}%
                                     {\normalfont\normalsize\bfseries}}
\newcommand*\l@subsubsubsection{\@dottedtocline{3}{10.0em}{4.1em}}
\newcommand*{\subsubsubsectionmark}[1]{}
\makeatother

\everymath{\color{black}}%make in-line maths symbols blue to read/check easily

\sloppy
\captionsetup[figure]{labelfont={small,bf},textfont={small,it},belowskip=-1pt,aboveskip=-9pt}
%space remove between caption, figure, and text
\captionsetup[table]{labelfont={small,bf},textfont={small,it},belowskip=-1pt,aboveskip=7pt}
\setlength{\floatsep}{5pt plus 2pt minus 2pt}
\setlength{\textfloatsep}{5pt plus 2pt minus 2pt}
\setlength{\intextsep}{10pt plus 2pt minus 2pt}

\thesislayout
\onehalfspacing

\begin{document}

\begin{titlepage}
\begin{tikzpicture}[remember picture, overlay]
  \draw[line width = 4pt] ($(current page.north west) + (0.4in,-0.5in)$) rectangle ($(current page.south east) + (-0.4in,0.5in)$);
  \draw[line width=1.5pt]
    ($ (current page.north west) + (0.45in,-0.55in) $)
    rectangle
    ($ (current page.south east) + (-0.45in,0.55in) $);
\end{tikzpicture}

\begin{center}
\LARGE \textbf{ĐẠI HỌC QUỐC GIA THÀNH PHỐ HỒ CHÍ MINH} \\
\vspace{0.15cm}
\LARGE \textbf{TRƯỜNG ĐẠI HỌC BÁCH KHOA} \\
\vspace{0.15cm}
\LARGE \textbf{KHOA KHOA HỌC VÀ KĨ THUẬT MÁY TÍNH}
\end{center}

\vspace{0.2cm}

\begin{figure}[h!]
\begin{center}

\includegraphics[width=7cm]{Images/hcmut.png}
\end{center}
\end{figure}

\begin{center}
\begin{tabular}{c}
\textbf{{\LARGE CO3084 - Mật mã học và Mã hóa Thông tin (TN)}} \\
\\
\textbf{{\LARGE Problem Set 1}} \\
\\
\parbox{0.5\textwidth}{%
    \centering
    \textbf{\LARGE Provable Security Foundations}
}
\end{tabular}
\end{center}

\begin{table}[h]
\begin{tabular}{rll}
\hspace{3.7cm} &   \textbf{\Large SVTH:  } {\Large Nguyễn Tấn Tài - 2212990}
\\
\\
\end{tabular}
\end{table}
\vspace{6cm}
\begin{center}
{\Large TP. Hồ Chí Minh, Tháng 9/2025 }
\end{center}
\end{titlepage}

\newpage

%%%%%%%%%%%%%%%%%INTRO%%%%%%%%%%%%%%%%%%%
\newpage

\tableofcontents

\newpage

\setcounter{page}{1}

%%%%%%%%%%%%%%%%%CONTENTS%%%%%%%%%%%%%%%%%%%
\section{Cryptographic Foundations}

\subsection{(10 điểm) Các khái niệm cơ bản}

\textbf{(a) (3 điểm) Định nghĩa ba mục tiêu an ninh chính của mật mã học bằng lời của bạn và đưa ra một ví dụ thực tế cho mỗi mục tiêu mà không được đề cập trực tiếp trong bài giảng.}

\begin{center}
\textbf{Trả lời:}
\end{center}

\textbf{Confidentiality (Bảo mật)}: Chỉ những bên được phép mới đọc được dữ liệu.
\begin{itemize}
    \item Nói cách khác, dù dữ liệu có bị chặn trên đường đi, nội dung vẫn không bị lộ.
    \item Ví dụ: Vault của trình quản lý mật khẩu (1Password/Bitwarden) được mã hóa end-to-end bằng khóa của người dùng thay vì lưu plain text trong cơ sở dữ liệu.
\end{itemize}

\textbf{Authentication (Xác thực)}: Xác minh danh tính thực sự của người gửi/người dùng.
\begin{itemize}
    \item Điều này giúp ta biết ``đúng người, đúng khóa'', tránh mạo danh.
    \item Ví dụ: Đăng nhập dùng khóa bảo mật phần cứng FIDO2/WebAuthn (YubiKey) để chứng minh danh tính mà không cần mật khẩu.
\end{itemize}

\textbf{Integrity (Toàn vẹn)}: Dữ liệu không bị thay đổi trong lưu trữ/truyền tải; nếu thay đổi sẽ bị phát hiện.
\begin{itemize}
    \item Tức là bất kỳ khác biệt nhỏ nào so với bản gốc cũng sẽ bị phát hiện một cách đáng tin cậy.
    \item Ví dụ: Kiểm tra chữ ký/tệp checksum SHA-256 kèm theo các bản cài đặt phần mềm (release checksum + signature) trước khi cài đặt.
\end{itemize}

\textbf{(b) (3 điểm) Giải thích nguyên lý Kerckhoff và lý do tại sao nguyên lý này vẫn là nền tảng của mật mã học hiện đại. Hãy đưa ra một ví dụ về một hệ thống bảo mật vi phạm nguyên lý này và mô tả những hậu quả tiềm ẩn.}

\begin{center}
\textbf{Trả lời:}
\end{center}

\textbf{Nguyên lý Kerckhoff}: Hệ mật mã vẫn phải an toàn ngay cả khi toàn bộ thuật toán/thiết kế/mã nguồn đều công khai; chỉ khóa là bí mật.

Có nghĩa là hệ mật mã chỉ cần bảo mật duy nhất khoá, các thứ khác (như là thuật toán, thiết kế, ...) thì có thể (hoặc là khuyến khích) công khai.

\textbf{Vì sao nền tảng}:

\begin{itemize}
    \item Tránh ``security through obscurity''; thuật toán lộ không làm sụp đổ hệ thống.
    \item Tập trung bảo vệ khóa (bí mật duy nhất), đơn giản hóa vận hành.
    \item Khuyến khích kiểm toán/đánh giá ngang hàng, tăng minh bạch và độ tin cậy, giảm nguy cơ backdoor.
\end{itemize}

\textbf{Ví dụ vi phạm}: Thiết bị IoT dùng ``mật mã tự chế'' đóng kín, dựa vào việc giấu thuật toán.
\begin{itemize}
    \item \textbf{Hậu quả}: Khi bị reverse-engineering rò thuật toán, kẻ tấn công có thể giải mã/giả mạo; không có đánh giá cộng đồng nên lỗ hổng tồn tại lâu; mất niềm tin người dùng.
\end{itemize}

Các tiêu chuẩn mở (AES, TLS) bền vững theo thời gian là vì hàng nghìn con mắt đã soi xét chúng, chứ không phải vì chúng bí mật, giấu kín.

\textbf{(c) (4 điểm) So sánh và đối chiếu giữa mật mã đối xứng và mật mã bất đối xứng:
    \begin{enumerate}
        \item Giải thích sự khác biệt cơ bản trong cách quản lý khóa của hai loại này.
        \item Đối với mỗi loại, hãy xác định các giả định toán học hoặc tính toán mà an ninh của chúng thường dựa vào.
        \item Mô tả một tình huống mà một loại sẽ rõ ràng là lựa chọn ưu việt hơn loại còn lại.
    \end{enumerate}
}

\begin{center}
\textbf{Trả lời:}
\end{center}

\begin{tabular}{|l|p{5cm}|p{5cm}|}
\hline
\textbf{Tiêu chí} & \textbf{Mật mã đối xứng} & \textbf{Mật mã bất đối xứng} \\
\hline
\textbf{Quản lý khóa} & 
1 khóa bí mật chung cho cả mã hóa và giải mã; cần phân phối khóa bí mật an toàn trước khi giao tiếp. & 
1 cặp khóa (công khai + riêng tư); khóa công khai chia sẻ rộng rãi, khóa riêng tư giữ bí mật; không cần kênh bí mật ban đầu. \\
\hline
\textbf{Cơ sở an ninh/giả định} & 
Dựa vào sức mạnh thuật toán và phân tích mật mã lâu dài (ví dụ: AES, HMAC-SHA-256). & 
Dựa vào các bài toán toán học khó (ví dụ: RSA dựa phân tích số nguyên lớn, DH/ECDH dựa logarit rời rạc). \\
\hline
\textbf{Tình huống ưu việt} & 
Mã hóa khối lượng dữ liệu lớn, yêu cầu hiệu năng cao (mã hóa đĩa, backup, streaming nội bộ sau khi đã có khóa phiên). & 
Trao đổi khóa ban đầu, xác thực/chữ ký số giữa các bên chưa có kênh bí mật (thiết lập TLS, ký cập nhật phần mềm). \\
\hline
\end{tabular}

\subsection{(10 điểm) Tính bí mật hoàn hảo (Perfect Secrecy)}

\textbf{(a) (3 điểm) Xét biến thể của one-time pad sử dụng phép AND bit ($\land$) thay vì XOR ($\oplus$):}

\begin{itemize}
    \item \textbf{Mã hóa:} $\text{Enc}(K, M) = K \land M$
    \item \textbf{Giải mã:} $\text{Dec}(K, C) = ?$
\end{itemize}

\begin{enumerate}
    \item \textbf{Scheme này có đúng không? Nếu đúng, hãy chỉ ra hàm giải mã. Nếu không, giải thích lý do.}
    
    \begin{center}
    \textbf{Trả lời:}
    \end{center}
    
    Scheme này \textbf{không đúng} (không đảm bảo tính đúng đắn).
    
    Phân tích chi tiết:
    \begin{itemize}
        \item Phép AND không khả nghịch vì có nhiều cặp $(K, M)$ khác nhau có thể tạo ra cùng ciphertext $C$.
        \item Ví dụ cụ thể: Nếu $C = 0100$, thì có thể có:
        \begin{itemize}
            \item $K = 0100, M = 0100 \rightarrow K \land M = 0100$
            \item $K = 0100, M = 0101 \rightarrow K \land M = 0100$
            \item $K = 0100, M = 0110 \rightarrow K \land M = 0100$
            \item $K = 0100, M = 0111 \rightarrow K \land M = 0100$
        \end{itemize}
        \item Khi giải mã với $K = 0100$ và $C = 0100$, không thể biết $M$ ban đầu là gì.
    \end{itemize}
    
    Tại sao AND không khả nghịch:
    \begin{itemize}
        \item Bit 0 trong $K$ ``che giấu'' thông tin về $M$ (luôn cho kết quả 0)
        \item Bit 1 trong $K$ ``tiết lộ'' thông tin về $M$ (kết quả = $M$)
        \item Thông tin bị mất không thể khôi phục được.
    \end{itemize}
    
    \item \textbf{Scheme này có đảm bảo bí mật hoàn hảo không? Giải thích.}
    
    \begin{center}
    \textbf{Trả lời:}
    \end{center}
    
    \textbf{Không đảm bảo bí mật hoàn hảo.}
    
    Phân tích chi tiết:
    \begin{itemize}
        \item Perfect Secrecy yêu cầu: $\Pr[\text{Enc}(K,M) = C] = 1/2^n$ với mọi $M$ và $C$
        \item Với phép AND, phân phối của ciphertext bị lệch nghiêm trọng
    \end{itemize}
    
    Ví dụ cụ thể với 2-bit:
    \begin{itemize}
        \item Nếu $M = 00$: chỉ có thể tạo ra $C = 00$ (dù $K$ là gì)
        \item Nếu $M = 01$: có thể tạo ra $C \in \{00, 01\}$ tùy theo $K$
        \item Nếu $M = 10$: có thể tạo ra $C \in \{00, 10\}$ tùy theo $K$
        \item Nếu $M = 11$: có thể tạo ra $C \in \{00, 01, 10, 11\}$ tùy theo $K$
    \end{itemize}
    
    Tấn công thông tin:
    \begin{itemize}
        \item Nếu attacker thấy $C = 11$, họ biết chắc $M = 11$ (vì chỉ có $M = 11$ mới tạo ra được $C = 11$)
        \item Nếu attacker thấy $C = 00$, họ biết $M$ có thể là bất kỳ giá trị nào (nhưng vẫn có thông tin)
        \item Điều này vi phạm định nghĩa Perfect Secrecy: attacker có thể phân biệt các $M$ khác nhau từ $C$.
    \end{itemize}
\end{enumerate}

\textbf{(b) (4 điểm) Xét biến thể one-time pad trên các chữ số thập phân (0-9):
    \begin{itemize}
        \item \textbf{Mã hóa:} $\text{Enc}(K, M) = (K + M) \bmod 10$
        \item \textbf{Giải mã:} $\text{Dec}(K, C) = (C - K) \bmod 10$
    \end{itemize}
}

Với $K, M, C \in \{0, 1, 2, \ldots, 9\}$

\begin{enumerate}
    \item \textbf{Chứng minh scheme này đúng}.
    
    \textbf{Chứng minh:}
    \begin{align}
        \text{Dec}(K, \text{Enc}(K, M)) &= ((K + M) \bmod 10 - K) \bmod 10 \\
        &= ((K + M - K) \bmod 10) \\
        &= (M \bmod 10) \\
        &= M \quad \text{(vì } M \in \{0,\ldots,9\}\text{)}
    \end{align}
    
    \textbf{Giải thích từng bước:}
    \begin{itemize}
        \item Bước 1: Thay thế $\text{Enc}(K,M) = (K + M) \bmod 10$
        \item Bước 2: Áp dụng tính chất phân phối của phép trừ trong modular arithmetic
        \item Bước 3: Đơn giản hóa $(K + M - K) = M$
        \item Bước 4: Vì $M \in \{0,1,\ldots,9\}$, nên $M \bmod 10 = M$
        \item Kết luận: $\text{Dec}(K, \text{Enc}(K,M)) = M$ với mọi $K$ và $M$
    \end{itemize}
    
    \item \textbf{Chứng minh scheme này đảm bảo bí mật hoàn hảo nếu $K$ chọn ngẫu nhiên đều.}
    
    \textbf{Chứng minh:}
    
    \textbf{Bước 1: Tính duy nhất của khóa}
    \begin{itemize}
        \item Với mỗi cặp $(M, C)$, tồn tại duy nhất một $K$ thỏa mãn: $C = (K + M) \bmod 10$
        \item Giải ra: $K = (C - M) \bmod 10$
    \end{itemize}
    
    \textbf{Bước 2: Phân phối đồng nhất của ciphertext}
    \begin{itemize}
        \item Nếu $K$ được chọn ngẫu nhiên đều từ $\{0,1,\ldots,9\}$, thì $\Pr[K = k] = 1/10$ với mọi $k$
        \item Với mỗi $M$ cố định, khi $K$ thay đổi từ 0 đến 9, ciphertext $C = (K + M) \bmod 10$ cũng chạy qua tất cả giá trị từ 0 đến 9
        \item Do đó: $\Pr[\text{Enc}(K,M) = C] = \Pr[K = (C-M) \bmod 10] = 1/10$ với mọi $M$ và $C$
    \end{itemize}
    
    \textbf{Bước 3: Kết luận Perfect Secrecy}
    \begin{itemize}
        \item Điều kiện Perfect Secrecy: $\Pr[\text{Enc}(K,M) = C] = 1/10$ với mọi $M$ và $C$
        \item Attacker không thể phân biệt các $M$ khác nhau từ $C$ vì phân phối của $C$ độc lập với $M$
    \end{itemize}
\end{enumerate}

Tổng kết: Cộng modulo là ``hoán vị dịch'' trên mỗi chữ số; nếu $K$ đồng nhất thì $K + M \bmod 10$ cũng đồng nhất, nên không lộ thông tin về $M$.

\textbf{(c) (3 điểm) Xét one-time pad với khóa ngắn bằng nửa độ dài thông điệp:}

\begin{itemize}
    \item $M = (M_1, M_2)$, với $|M_1| = |M_2| = |K|$
    \item $\text{Enc}(K, M) = (K \oplus M_1, K \oplus M_2)$
\end{itemize}

\textbf{Yêu cầu:} Đưa ra tấn công cụ thể phá vỡ tính bảo mật, chỉ rõ attacker có thể rút ra thông tin gì từ ciphertext.

\vspace{0.5em}
\textbf{Trả lời:}

\begin{enumerate}
    \item \textbf{Thu thập ciphertext:}
    \begin{itemize}
        \item Ciphertext thu được là $(C_1, C_2) = (K \oplus M_1, K \oplus M_2)$.
    \end{itemize}

    \item \textbf{Tấn công loại bỏ khóa:}
    \begin{itemize}
        \item Attacker tính $C_1 \oplus C_2$:
        \begin{align*}
            C_1 \oplus C_2 &= (K \oplus M_1) \oplus (K \oplus M_2) \\
            &= K \oplus M_1 \oplus K \oplus M_2 \\
            &= (K \oplus K) \oplus (M_1 \oplus M_2) \\
            &= 0 \oplus (M_1 \oplus M_2) \\
            &= M_1 \oplus M_2
        \end{align*}
        \item Như vậy, attacker biết được giá trị $M_1 \oplus M_2$.
    \end{itemize}

    \item \textbf{Phân tích thông tin bị lộ:}
    \begin{itemize}
        \item Giá trị $M_1 \oplus M_2$ tiết lộ mối quan hệ giữa hai nửa thông điệp.
        \item Ví dụ: Nếu $M_1 \oplus M_2 = 1010$, attacker có thể thử các giá trị $M_1$ để suy ra $M_2 = M_1 \oplus 1010$.
        \item Điều này giúp attacker loại trừ nhiều khả năng về thông điệp gốc, tức là ciphertext đã tiết lộ thông tin về plaintext.
    \end{itemize}

    \item \textbf{Giải thích vi phạm Perfect Secrecy:}
    \begin{itemize}
        \item Theo định nghĩa, Perfect Secrecy yêu cầu ciphertext không tiết lộ bất kỳ thông tin nào về plaintext.
        \item Tuy nhiên, ở đây attacker luôn biết $M_1 \oplus M_2$ từ ciphertext, nên scheme này không đạt Perfect Secrecy.
    \end{itemize}

    \item Việc sử dụng cùng một khóa $K$ cho hai phần thông điệp khiến khi XOR hai bản mã, khóa bị triệt tiêu, tương tự như bài toán two-time pad.

\end{enumerate}

% \newpage

% \input{Contents/2.Tong_quan_cong_viec}

% \newpage

% \section{Cơ sở lý thuyết}
\subsection{ Các yếu tố ảnh hưởng đến sản lượng nông sản}

Giá nông sản là một biến số nhạy cảm, chịu ảnh hưởng bởi nhiều yếu tố phức tạp từ cả phía cung và cầu. Việc hiểu rõ các yếu tố này không chỉ giúp phân tích chính xác mà còn là nền tảng để xây dựng các mô hình dự đoán hiệu quả. Một số yếu tố chủ yếu bao gồm:

\begin{itemize}
    \item \textbf{Thời tiết và khí hậu:} Nhiệt độ, lượng mưa, hạn hán, lũ lụt và các hiện tượng thời tiết cực đoan ảnh hưởng trực tiếp đến năng suất và chất lượng cây trồng.
    
    \item \textbf{Dịch bệnh và sâu bệnh:} Sự bùng phát của sâu bệnh làm giảm sản lượng và chất lượng nông sản, từ đó đẩy giá thành lên cao.

    \item \textbf{Chi phí đầu vào:} Bao gồm giá giống cây trồng, phân bón, nhân công và đặc biệt là \textbf{thuốc trừ sâu}. Việc sử dụng thuốc trừ sâu không chỉ ảnh hưởng đến chi phí sản xuất mà còn tác động đến chất lượng sản phẩm, từ đó ảnh hưởng đến giá cả.

    \item \textbf{Chính sách nhà nước:} Các chính sách về thuế, trợ cấp, khuyến khích sản xuất hoặc hạn chế xuất khẩu cũng ảnh hưởng lớn đến giá nông sản.

    \item \textbf{Quan hệ cung - cầu thị trường:} Khi sản lượng vượt cầu, giá sẽ giảm; ngược lại khi nguồn cung hạn chế, giá có xu hướng tăng.

    \item \textbf{Tác động từ thị trường quốc tế:} Biến động giá cả toàn cầu, tỉ giá hối đoái và tình hình xuất nhập khẩu đều tác động đến giá nông sản trong nước.
\end{itemize}

Những yếu tố này thường có mối quan hệ phức tạp và phi tuyến tính với giá nông sản, khiến cho việc phân tích và dự đoán trở thành một bài toán đầy thách thức nhưng cũng rất cần thiết và có tính ứng dụng cao trong thực tiễn.


\subsection{ Tổng quan về học máy và phân lớp dữ liệu}

Học máy (\textit{Machine Learning}) là một nhánh của trí tuệ nhân tạo (AI), cho phép máy tính học từ dữ liệu và cải thiện hiệu suất dự đoán mà không cần lập trình một cách cụ thể. Trong bối cảnh dự đoán giá nông sản, học máy đóng vai trò quan trọng trong việc xây dựng các mô hình có khả năng khai thác các mối quan hệ phức tạp giữa nhiều yếu tố ảnh hưởng khác nhau.

Học máy được chia thành nhiều loại, trong đó phổ biến nhất là:
\begin{itemize}
    \item \textbf{Học có giám sát (Supervised Learning):} Mô hình được huấn luyện trên tập dữ liệu có nhãn, tức là mỗi mẫu dữ liệu đều đi kèm với kết quả đầu ra mong muốn. Đây là phương pháp chính được sử dụng trong bài toán dự đoán giá nông sản.
    
    \item \textbf{Học không giám sát (Unsupervised Learning):} Tập trung vào việc tìm kiếm cấu trúc ẩn trong dữ liệu không có nhãn (ví dụ: phân cụm).
    
    \item \textbf{Học tăng cường (Reinforcement Learning):} Mô hình học thông qua tương tác với môi trường và tối ưu hóa phần thưởng.
\end{itemize}

Trong khuôn khổ đồ án này, bài toán chính được xác định là một bài toán \textbf{hồi quy}, trong đó đầu ra cần dự đoán là một giá trị liên tục (giá nông sản). Tuy nhiên, cũng có thể tiếp cận dưới dạng \textbf{phân lớp} nếu ta chia giá thành các mức giá cụ thể như “thấp”, “trung bình” và “cao”.

Các thuật toán học máy thường được sử dụng trong phân tích và dự đoán dữ liệu gồm có: hồi quy tuyến tính, cây quyết định, xác suất Bayes (Naive Bayes), mạng nơ-ron nhân tạo và K-nearest neighbors (KNN). Mỗi thuật toán có đặc điểm riêng phù hợp với từng kiểu dữ liệu và yêu cầu dự đoán khác nhau, sẽ được trình bày cụ thể ở phần tiếp theo.


\subsection{ Các thuật toán học máy}

Trong quá trình xây dựng mô hình dự đoán giá nông sản, việc lựa chọn thuật toán phù hợp là yếu tố then chốt để đảm bảo độ chính xác và khả năng tổng quát hóa của mô hình. Dưới đây là các thuật toán phổ biến được sử dụng trong bài toán dự đoán giá trị liên tục hoặc phân lớp giá trị.

\subsubsection{ Hồi quy tuyến tính (Linear Regression)}

Hồi quy tuyến tính là một trong những mô hình cơ bản nhất trong học máy, được sử dụng rộng rãi trong các bài toán dự đoán. Mục tiêu của hồi quy tuyến tính là tìm ra một hàm tuyến tính \( y = w_0 + w_1x_1 + w_2x_2 + \ldots + w_nx_n \) sao cho sai số giữa giá trị dự đoán và giá trị thực tế là nhỏ nhất.

\begin{itemize}
    \item \textbf{Ưu điểm:} Dễ hiểu, dễ triển khai, thời gian huấn luyện nhanh.
    \item \textbf{Nhược điểm:} Chỉ phù hợp với các mối quan hệ tuyến tính giữa các biến; độ chính xác không cao khi dữ liệu có quan hệ phi tuyến.
\end{itemize}

\subsubsection{ Cây quyết định (Decision Tree)}

Cây quyết định là mô hình dựa trên cấu trúc cây phân nhánh, trong đó mỗi nút là một điều kiện phân chia dữ liệu dựa trên giá trị của một thuộc tính, và các nhánh dẫn đến các quyết định hoặc dự đoán.

\begin{itemize}
    \item \textbf{Ưu điểm:} Dễ giải thích, không cần chuẩn hóa dữ liệu, xử lý tốt dữ liệu hỗn hợp.
    \item \textbf{Nhược điểm:} Dễ bị overfitting nếu cây quá sâu, kém ổn định với dữ liệu nhiễu.
\end{itemize}


\subsubsection{ Mạng nơ-ron nhân tạo (Artificial Neural Network - ANN)}

Mạng nơ-ron nhân tạo là mô hình học sâu được lấy cảm hứng từ cách hoạt động của bộ não con người. Một mạng nơ-ron bao gồm các lớp (layer) gồm nhiều “nơ-ron” (neuron), được kết nối với nhau thông qua các trọng số. Dữ liệu đầu vào được truyền qua từng lớp (input, hidden, output), trải qua các hàm kích hoạt và được huấn luyện bằng cách cập nhật trọng số để giảm sai số dự đoán.

\begin{itemize}
    \item \textbf{Ưu điểm:} Khả năng mô hình hóa các mối quan hệ phi tuyến phức tạp, phù hợp với dữ liệu lớn và nhiều chiều.
    \item \textbf{Nhược điểm:} Cần nhiều tài nguyên tính toán, dễ overfitting nếu không điều chỉnh tốt; khó giải thích kết quả so với các mô hình truyền thống.
\end{itemize}

Trong bài toán dự đoán giá nông sản, ANN đặc biệt hiệu quả khi có dữ liệu lớn, với nhiều đặc trưng đầu vào như thời tiết, phân bón, thuốc trừ sâu, v.v.

\subsubsection{ K-Nearest Neighbors (KNN)}

KNN là thuật toán đơn giản nhưng mạnh mẽ, dựa trên nguyên lý “gần mực thì đen, gần đèn thì sáng”. Khi cần dự đoán một điểm mới, KNN sẽ tìm \( k \) điểm dữ liệu gần nhất trong không gian đặc trưng, sau đó lấy trung bình (đối với bài toán hồi quy) hoặc đa số (đối với phân lớp) làm kết quả.

\begin{itemize}
    \item \textbf{Ưu điểm:} Dễ triển khai, không giả định phân phối dữ liệu.
    \item \textbf{Nhược điểm:} Tốc độ chậm khi dữ liệu lớn, nhạy cảm với nhiễu và khoảng cách.
\end{itemize}

Việc lựa chọn giá trị \( k \) phù hợp là yếu tố then chốt quyết định hiệu quả của mô hình KNN.




% \newpage

% \input{Contents/4.Tien_xu_ly_du_lieu}

% \newpage

% \input{Contents/5.TrucQuanHoaDuLieu}

% \newpage

% \section{Xây dựng mô hình dự đoán}
\subsection{Các đại lượng đánh giá mô hình}
\label{subsec:evaluation_metrics_explanation}

Để đánh giá chất lượng và độ chính xác của các mô hình trong việc dự đoán sản lượng cây trồng, nhóm em sử dụng một tập hợp các đại lượng (metrics) đo lường phổ biến. Các đại lượng này giúp định lượng mức độ lỗi của mô hình cũng như khả năng giải thích sự biến thiên trong dữ liệu của mô hình. Dưới đây là giải thích chi tiết về từng đại lượng:

Trong các công thức dưới đây, chúng ta ký hiệu:
\begin{itemize}
    \item $n$: tổng số quan sát trong tập dữ liệu đánh giá (ví dụ: tập kiểm tra).
    \item $y_i$: giá trị thực tế của quan sát thứ $i$.
    \item $\hat{y}_i$: giá trị dự đoán bởi mô hình cho quan sát thứ $i$.
    \item $\bar{y}$: giá trị trung bình của các giá trị thực tế $y_i$, tức là $\bar{y} = \frac{1}{n} \sum_{i=1}^{n} y_i$.
\end{itemize}

\subsubsection{Mean Absolute Error (MAE) – Sai số Tuyệt đối trung bình}
\label{subsubsec:mae_explanation}
MAE đo lường độ lớn trung bình của các sai số trong một tập hợp các dự đoán, mà không xem xét đến chiều hướng của chúng. Đây là trung bình cộng của các chênh lệch tuyệt đối giữa giá trị dự đoán và giá trị thực tế.
\begin{equation}
    \text{MAE} = \frac{1}{n} \sum_{i=1}^{n} |y_i - \hat{y}_i|
    \label{eq:mae}
\end{equation}
\textbf{Ý nghĩa và diễn giải:}
\begin{itemize}
    \item MAE cho chúng ta biết, trung bình, dự đoán của mô hình lệch khỏi giá trị thực tế bao nhiêu đơn vị.
    \item Đơn vị của MAE giống với đơn vị của biến mục tiêu (ví dụ: nếu sản lượng tính bằng tấn/ha, MAE cũng sẽ có đơn vị là tấn/ha).
    \item Giá trị MAE càng gần 0, mô hình dự đoán càng chính xác.
    \item MAE không bị ảnh hưởng nhiều bởi các giá trị ngoại lệ (outliers) lớn như MSE vì nó không bình phương sai số.
\end{itemize}

\subsubsection{Mean Squared Error (MSE) – Sai số Bình phương trung bình}
\label{subsubsec:mse_explanation}
MSE là trung bình của bình phương các sai số, tức là chênh lệch giữa giá trị ước lượng và giá trị thực tế.
\begin{equation}
    \text{MSE} = \frac{1}{n} \sum_{i=1}^{n} (y_i - \hat{y}_i)^2
    \label{eq:mse}
\end{equation}
\textbf{Ý nghĩa và diễn giải:}
\begin{itemize}
    \item MSE nhấn mạnh các lỗi lớn hơn bằng cách bình phương chúng trước khi lấy trung bình. Do đó, một mô hình có MSE cao có thể là do có một vài dự đoán với sai số rất lớn.
    \item Đơn vị của MSE là bình phương của đơn vị biến mục tiêu (ví dụ: (tấn/ha)²), điều này làm cho việc diễn giải trực tiếp giá trị MSE trở nên khó khăn hơn so với MAE hoặc RMSE.
    \item Giá trị MSE càng gần 0, mô hình càng tốt.
    \item MSE thường được sử dụng trong các hàm mất mát (loss function) của nhiều thuật toán học máy do tính chất toán học thuận lợi của nó (ví dụ: khả vi).
\end{itemize}

\subsubsection{Root Mean Squared Error (RMSE) – Căn bậc hai Sai số Bình phương trung bình}
\label{subsubsec:rmse_explanation}
RMSE là căn bậc hai của MSE. Nó là một trong những đại lượng phổ biến nhất để đo lường sự khác biệt giữa các giá trị được dự đoán bởi một mô hình và các giá trị thực tế.
\begin{equation}
    \text{RMSE} = \sqrt{\text{MSE}} = \sqrt{\frac{1}{n} \sum_{i=1}^{n} (y_i - \hat{y}_i)^2}
    \label{eq:rmse}
\end{equation}
\textbf{Ý nghĩa và diễn giải:}
\begin{itemize}
    \item RMSE có cùng đơn vị với biến mục tiêu (giống như MAE), giúp việc diễn giải trở nên trực quan hơn so với MSE. Ví dụ, nếu RMSE là 0.5 tấn/ha, điều này có nghĩa là độ lệch chuẩn của các phần dư (sai số dự đoán) là 0.5 tấn/ha.
    \item Tương tự MSE, RMSE cũng nhạy cảm với các giá trị ngoại lệ do việc bình phương sai số. Các lỗi lớn sẽ làm tăng RMSE một cách đáng kể.
    \item Giá trị RMSE càng gần 0, mô hình dự đoán càng chính xác.
    \item Khi so sánh hai mô hình, mô hình có RMSE thấp hơn thường được coi là tốt hơn (giả sử các yếu tố khác là tương đương).
\end{itemize}

\subsubsection{R-squared (R²) – Hệ số xác định}
\label{subsubsec:r_squared_explanation}
R-squared, hay còn gọi là hệ số xác định, là một đại lượng thống kê đo lường tỷ lệ phần trăm phương sai (variance) của biến phụ thuộc (biến mục tiêu) mà mô hình hồi quy có thể giải thích được. Nói cách khác, nó cho biết mức độ phù hợp của mô hình với dữ liệu quan sát được.
\begin{equation}
    R^2 = 1 - \frac{SS_{\text{res}}}{SS_{\text{tot}}} = 1 - \frac{\sum_{i=1}^{n} (y_i - \hat{y}_i)^2}{\sum_{i=1}^{n} (y_i - \bar{y})^2}
    \label{eq:r_squared}
\end{equation}
Trong đó:
\begin{itemize}
    \item $SS_{\text{res}} = \sum_{i=1}^{n} (y_i - \hat{y}_i)^2$ là tổng bình phương các phần dư (Residual Sum of Squares), đo lường phần phương sai không giải thích được bởi mô hình.
    \item $SS_{\text{tot}} = \sum_{i=1}^{n} (y_i - \bar{y})^2$ là tổng bình phương toàn phần (Total Sum of Squares), đo lường tổng phương sai của biến mục tiêu.
\end{itemize}
\textbf{Ý nghĩa và diễn giải:}
\begin{itemize}
    \item $R^2$ có giá trị nằm trong khoảng từ $-\infty$ đến $1$. Tuy nhiên, trong hầu hết các trường hợp thực tế với các mô hình hợp lý, $R^2$ sẽ nằm trong khoảng từ $0$ đến $1$.
    \item $R^2 = 1$: Mô hình giải thích được toàn bộ sự biến thiên trong dữ liệu, tức là các giá trị dự đoán hoàn toàn khớp với giá trị thực tế. Đây là trường hợp lý tưởng nhưng hiếm khi xảy ra.
    \item $R^2 = 0$: Mô hình không giải thích được chút nào sự biến thiên trong dữ liệu. Mô hình dự đoán không tốt hơn việc chỉ sử dụng giá trị trung bình $\bar{y}$ để dự đoán cho tất cả các quan sát.
    \item $R^2 > 0$: Ví dụ, $R^2 = 0.75$ có nghĩa là $75\%$ sự biến thiên của biến mục tiêu được giải thích bởi các biến độc lập trong mô hình. $25\%$ còn lại là do các yếu tố khác không được mô hình nắm bắt.
    \item $R^2 < 0$: Điều này có thể xảy ra nếu mô hình dự đoán tệ hơn cả việc chỉ dùng giá trị trung bình $\bar{y}$. Thường gặp khi mô hình được chọn không phù hợp với dữ liệu.
    \item Giá trị $R^2$ càng gần $1$ thì mô hình càng phù hợp với dữ liệu. Tuy nhiên, một $R^2$ cao không phải lúc nào cũng đồng nghĩa với một mô hình tốt. Việc thêm quá nhiều biến vào mô hình (ngay cả những biến không liên quan) có thể làm tăng $R^2$ một cách giả tạo. Do đó, $R^2$ nên được xem xét cùng với các chỉ số khác và bối cảnh của bài toán. Adjusted R-squared (R² điều chỉnh) là một biến thể thường được sử dụng để khắc phục nhược điểm này bằng cách phạt việc thêm các biến không cần thiết.
\end{itemize}

Việc hiểu rõ ý nghĩa của từng đại lượng đánh giá này giúp chúng ta có cái nhìn toàn diện hơn về hiệu suất của các mô hình hồi quy và đưa ra những quyết định tốt hơn trong việc lựa chọn và cải tiến mô hình.






\subsection{Mô hình Hồi quy tuyến tính}
\label{subsec:linear_regression_results}

Sau khi xây dựng và huấn luyện mô hình Hồi quy Tuyến tính sử dụng các đặc trưng về thời tiết, nông nghiệp và các yếu tố khác (bao gồm các biến phân loại đã được mã hóa bằng One-Hot Encoding) để dự đoán sản lượng cây trồng (`Crop\_Yield\_MT\_per\_HA`), nhóm em đã thu được các kết quả đánh giá và hệ số như sau. 

\subsubsection{Các chỉ số đánh giá hiệu suất mô hình}
\label{subsubsec:lr_performance_metrics}

Hiệu suất của mô hình Hồi quy Tuyến tính trên tập kiểm tra được đánh giá thông qua các chỉ số sau:
\begin{itemize}
    \item \textbf{Mean Absolute Error (MAE):} $0.5466$
    \item \textbf{Mean Squared Error (MSE):} $0.4648$
    \item \textbf{Root Mean Squared Error (RMSE):} $0.6818$
    \item \textbf{R-squared (R²):} $0.5596$
\end{itemize}
Hệ số chặn (Intercept) của mô hình là $-0.6208$.

Giá trị $R^2 = 0.5596$ cho thấy mô hình Hồi quy Tuyến tính của nhóm em giải thích được khoảng $55.96\%$ sự biến thiên của sản lượng cây trồng. Đây là một mức độ giải thích ở mức khá, chỉ ra rằng các đặc trưng được lựa chọn có vai trò nhất định trong việc dự đoán sản lượng. Tuy nhiên, cũng có nghĩa là còn khoảng $44\%$ sự biến thiên của sản lượng chưa được mô hình nắm bắt, có thể do các yếu tố khác không được đưa vào mô hình hoặc do mối quan hệ phi tuyến tính phức tạp hơn.

Các chỉ số sai số MAE ($0.5466$) và RMSE ($0.6818$) cho biết mức độ lỗi trung bình của dự đoán so với giá trị thực tế (đơn vị là tấn/ha, giả sử sản lượng được đo bằng đơn vị này). Ví dụ, RMSE là $0.6818$ tấn/ha ngụ ý rằng, trung bình, các dự đoán của mô hình có thể lệch khoảng $0.68$ tấn/ha so với sản lượng thực tế. Mức độ sai số này cần được xem xét trong bối cảnh giá trị trung bình và độ biến động chung của sản lượng cây trồng trong bộ dữ liệu.

\subsubsection{Phân tích các hệ số hồi quy}
\label{subsubsec:lr_coefficients}

Các hệ số hồi quy ước lượng mức độ ảnh hưởng của mỗi đặc trưng đến sản lượng cây trồng, khi các yếu tố khác được giữ không đổi. Dưới đây là một số hệ số nổi bật từ mô hình (đã được sắp xếp theo giá trị tuyệt đối giảm dần từ file \texttt{regression\_coefficients.csv}):

\begin{itemize}
    \item \texttt{Crop\_Type\_Rice}: $2.2917$
    \item \texttt{Crop\_Type\_Wheat}: $-1.3206$
    \item \texttt{Crop\_Type\_Soybeans}: $-1.0568$
    \item \texttt{Fertilizer\_Use\_KG\_per\_HA}: $0.8000$
    \item \texttt{Adaptation\_Strategies\_Water Management}: $-0.7121$
\end{itemize}

Từ các hệ số này, có thể rút ra một số nhận xét sơ bộ:
\begin{itemize}
    \item Việc trồng lúa (\texttt{Crop\_Type\_Rice}) có hệ số dương lớn nhất ($2.2917$), cho thấy loại cây này có sản lượng dự kiến cao hơn đáng kể so với loại cây tham chiếu (sau khi One-Hot Encoding). Ngược lại, lúa mì (\texttt{Crop\_Type\_Wheat}) và đậu nành (\texttt{Crop\_Type\_Soybeans}) có hệ số âm, ngụ ý sản lượng dự kiến thấp hơn.
    \item Lượng phân bón sử dụng (\texttt{Fertilizer\_Use\_KG\_per\_HA}) có hệ số dương ($0.8000$), điều này phù hợp với kỳ vọng rằng việc tăng cường sử dụng phân bón (trong một chừng mực nhất định) sẽ giúp tăng sản lượng.
    \item Chiến lược thích ứng về quản lý nước (\texttt{Adaptation\_Strategies\_Water Management}) có hệ số âm ($-0.7121$). Điều này có vẻ phản trực giác và cần được phân tích sâu hơn. Có thể các khu vực áp dụng chiến lược này vốn dĩ đã gặp vấn đề nghiêm trọng về nước, hoặc chiến lược này chưa thực sự hiệu quả trong bối cảnh dữ liệu thu thập được, hoặc nó tương quan với các yếu tố khác làm giảm sản lượng.
\end{itemize}
Cần lưu ý rằng đây là các diễn giải ban đầu và mối quan hệ thực tế có thể phức tạp hơn do tương tác giữa các biến. Toàn bộ danh sách các hệ số được lưu trong file \texttt{regression\_coefficients.csv} tại thư mục kết quả.

\subsubsection{Trực quan hóa kết quả}
\label{subsubsec:lr_visualization}

Để đánh giá trực quan hơn về hiệu suất của mô hình, nhóm em xem xét hai biểu đồ chính: biểu đồ so sánh giá trị sản lượng thực tế và dự đoán (Hình \ref{fig:lr_actual_vs_predicted}), và biểu đồ phân phối của phần dư (Hình \ref{fig:lr_residuals_distribution}).

\begin{figure}[htb]
    \centering
    \includegraphics[width=0.9\textwidth]{Images/actual_vs_predicted_plot.png} 
    \caption{So sánh giá trị sản lượng thực tế và giá trị dự đoán bởi mô hình Hồi quy Tuyến tính trên tập kiểm tra.}
    \vspace{10pt}
    \label{fig:lr_actual_vs_predicted}
\end{figure}
\FloatBarrier
Hình \ref{fig:lr_actual_vs_predicted} cho thấy các điểm dữ liệu (mỗi điểm tương ứng với một quan sát trong tập kiểm tra) phân bố quanh đường chéo $y=x$. Nếu các điểm nằm hoàn toàn trên đường này, mô hình sẽ dự đoán chính xác tuyệt đối. Trong trường hợp này, các điểm có xu hướng bám theo đường chéo, nhưng cũng có sự phân tán nhất định, phản ánh các sai số MAE và RMSE đã tính toán ở trên. Điều này cho thấy mô hình nắm bắt được xu hướng chung nhưng vẫn còn những dự đoán chưa hoàn toàn chính xác.

\begin{figure}[htb]
    \centering
    \includegraphics[width=0.9\textwidth]{Images/residuals_distribution_plot.png} 
    \vspace{10pt}
    \caption{Phân phối của phần dư (Actual - Predicted) từ mô hình Hồi quy Tuyến tính trên tập kiểm tra.}
    \label{fig:lr_residuals_distribution}
\end{figure}
\FloatBarrier
Hình \ref{fig:lr_residuals_distribution} thể hiện phân phối của phần dư (sai số giữa giá trị thực tế và giá trị dự đoán). Một mô hình tốt thường có phần dư phân phối chuẩn với giá trị trung bình gần bằng 0 và không có hình mẫu rõ ràng khi vẽ theo giá trị dự đoán (điều này không được hiển thị ở đây nhưng là một phần của phân tích phần dư sâu hơn). Trong biểu đồ này, phần dư có vẻ tập trung quanh 0 và có hình dạng gần giống phân phối chuông, điều này là một dấu hiệu tích cực, cho thấy mô hình không có thiên lệch hệ thống (systematic bias) rõ rệt. Tuy nhiên, sự tồn tại của một số phần dư có giá trị lớn (ở hai phía của đuôi phân phối) cũng chỉ ra các trường hợp mà mô hình dự đoán kém chính xác.

\subsubsection{Tóm tắt và Nhận xét}
Mô hình Hồi quy Tuyến tính đã cung cấp một cái nhìn ban đầu về mối quan hệ giữa các yếu tố đầu vào và sản lượng nông nghiệp. Với $R^2 \approx 0.56$, mô hình có khả năng giải thích một phần đáng kể sự biến động của sản lượng. Các hệ số cung cấp thông tin về hướng và cường độ tương đối của các yếu tố. Tuy nhiên, mức độ sai số và phần $R^2$ còn lại cho thấy tiềm năng cải thiện bằng cách sử dụng các mô hình phức tạp hơn, kỹ thuật đặc trưng tiên tiến hơn, hoặc thu thập thêm dữ liệu/đặc trưng có liên quan.





\subsection{ Mô hình Cây Quyết định (Decision Tree Regressor)}
\label{subsec:decision_tree_results}

Tiếp theo mô hình Hồi quy Tuyến tính, nhóm em đã xây dựng và đánh giá một mô hình Cây Quyết định (Decision Tree Regressor) để dự đoán sản lượng cây trồng (`Crop\_Yield\_MT\_per\_HA`). Mô hình này sử dụng các đặc trưng tương tự như mô hình hồi quy, bao gồm các biến đã được mã hóa One-Hot.

\subsubsection{Các chỉ số đánh giá hiệu suất mô hình}
\label{subsubsec:dt_performance_metrics}

Hiệu suất của mô hình Cây Quyết định trên tập kiểm tra được đánh giá như sau:
\begin{itemize}
    \item \textbf{Mean Absolute Error (MAE):} $0.6705$
    \item \textbf{Mean Squared Error (MSE):} $0.7727$
    \item \textbf{Root Mean Squared Error (RMSE):} $0.8791$
    \item \textbf{R-squared (R²):} $0.2679$
\end{itemize}

Giá trị $R^2 = 0.2679$ cho thấy mô hình Cây Quyết định (với các tham số mặc định) giải thích được khoảng $26.79\%$ sự biến thiên của sản lượng cây trồng. Con số này thấp hơn đáng kể so với mô hình Hồi quy Tuyến tính ($R^2 \approx 0.56$). Điều này gợi ý rằng mô hình Cây Quyết định cơ bản có thể chưa nắm bắt tốt các mối quan hệ trong dữ liệu, hoặc có thể đang bị hiện tượng overfitting (học quá khớp trên tập huấn luyện) hoặc underfitting nếu cây quá nông.

Các chỉ số sai số MAE ($0.6705$) và RMSE ($0.8791$) cũng cao hơn so với mô hình Hồi quy Tuyến tính, cho thấy dự đoán của mô hình Cây Quyết định này, trung bình, lệch nhiều hơn so với giá trị sản lượng thực tế.

\subsubsection{Phân tích độ quan trọng của các đặc trưng (Feature Importances)}
\label{subsubsec:dt_feature_importances}

Khác với Hồi quy Tuyến tính sử dụng hệ số, Cây Quyết định đánh giá tầm quan trọng của các đặc trưng dựa trên mức độ chúng đóng góp vào việc giảm thiểu sai số (hoặc tạp chất) tại các nút chia của cây. Dưới đây là 5 đặc trưng có độ quan trọng cao nhất theo mô hình (từ file \texttt{dt\_feature\_importances.csv}):

\begin{enumerate}
    \item \texttt{Fertilizer\_Use\_KG\_per\_HA}: $0.5360$
    \item \texttt{Crop\_Type\_Rice}: $0.1608$
    \item \texttt{Year}: $0.0616$
    \item \texttt{Pesticide\_Use\_KG\_per\_HA}: $0.0400$
    \item \texttt{Soil\_Health\_Index}: $0.0331$
\end{enumerate}

Nhận xét từ các đặc trưng quan trọng nhất:
\begin{itemize}
    \item Lượng phân bón sử dụng (\texttt{Fertilizer\_Use\_KG\_per\_HA}) là đặc trưng có ảnh hưởng lớn nhất đến các quyết định của cây, chiếm tới hơn $53\%$ tổng độ quan trọng. Điều này nhấn mạnh vai trò then chốt của phân bón trong mô hình dự đoán sản lượng này.
    \item Việc trồng lúa (\texttt{Crop\_Type\_Rice}) cũng là một yếu tố quan trọng ($16.08\%$).
    \item Năm (\texttt{Year}), lượng thuốc trừ sâu (\texttt{Pesticide\_Use\_KG\_per\_HA}), và chỉ số sức khỏe đất (\texttt{Soil\_Health\_Index}) cũng có những đóng góp đáng kể, mặc dù ở mức độ thấp hơn.
\end{itemize}
Việc các đặc trưng liên quan đến loại cây trồng và yếu tố đầu vào nông nghiệp (phân bón, thuốc trừ sâu) có độ quan trọng cao là điều dễ hiểu. Độ quan trọng của `Year` có thể phản ánh xu hướng thay đổi sản lượng theo thời gian do các yếu tố không được đo lường trực tiếp khác (ví dụ: cải tiến công nghệ, chính sách nông nghiệp).

\subsubsection{Trực quan hóa kết quả và cấu trúc cây}
\label{subsubsec:dt_visualization}

Các biểu đồ sau giúp đánh giá trực quan hiệu suất và cấu trúc của mô hình Cây Quyết định.

\begin{figure}[htb]
    \centering
    \includegraphics[width=0.8\textwidth]{Images/dt_actual_vs_predicted_plot.png} % ĐIỀU CHỈNH ĐƯỜNG DẪN
    \caption{So sánh giá trị sản lượng thực tế và giá trị dự đoán bởi mô hình Cây Quyết định trên tập kiểm tra.}
    \label{fig:dt_actual_vs_predicted}
\end{figure}

Hình \ref{fig:dt_actual_vs_predicted} cho thấy sự phân tán của các điểm dự đoán so với giá trị thực tế. So với mô hình Hồi quy Tuyến tính, các điểm có vẻ phân tán rộng hơn so với đường chéo $y=x$, điều này cũng phù hợp với giá trị $R^2$ thấp hơn và các chỉ số lỗi cao hơn. Có thể thấy một số cụm điểm dự đoán nằm trên các đường ngang, đây là đặc điểm của Cây Quyết định khi các mẫu rơi vào cùng một nút lá sẽ có cùng một giá trị dự đoán.

\begin{figure}[htb]
    \centering
    \includegraphics[width=0.8\textwidth]{Images/dt_residuals_distribution_plot.png} 
    \caption{Phân phối của phần dư (Actual - Predicted) từ mô hình Cây Quyết định trên tập kiểm tra.}
    \label{fig:dt_residuals_distribution}
\end{figure}

Phân phối phần dư (Hình \ref{fig:dt_residuals_distribution}) của mô hình Cây Quyết định cũng tập trung quanh 0, nhưng có vẻ rộng hơn và không đối xứng bằng so với mô hình Hồi quy Tuyến tính. Điều này cũng củng cố nhận định rằng mô hình Cây Quyết định (với tham số mặc định) có hiệu suất kém hơn trong trường hợp này.

\begin{figure}[htb]
    \centering
    \includegraphics[width=0.9\textwidth]{Images/dt_tree_visualization.png} 
    \caption{Trực quan hóa 3 tầng đầu tiên của mô hình Cây Quyết định. Các nút hiển thị điều kiện chia, giá trị mse, số lượng mẫu (samples) và giá trị dự đoán trung bình (value) tại nút đó.}
    \label{fig:dt_tree_visualization}
\end{figure}

Hình \ref{fig:dt_tree_visualization} minh họa cấu trúc của 3 tầng đầu tiên trong cây quyết định đã huấn luyện. Tại mỗi nút, cây đưa ra một quyết định dựa trên một đặc trưng và một ngưỡng giá trị để chia dữ liệu. Ví dụ, ở nút gốc (tầng trên cùng), cây có thể chia dựa trên \texttt{Fertilizer\_Use\_KG\_per\_HA}. Các mẫu dữ liệu sau đó sẽ đi theo các nhánh tương ứng dựa trên việc chúng thỏa mãn điều kiện chia hay không. Giá trị `value` tại mỗi nút lá (hoặc nút trung gian) là giá trị sản lượng trung bình của các mẫu thuộc về nút đó. Biểu đồ này giúp hiểu rõ hơn cách mô hình đưa ra dự đoán, dù chỉ là một phần nhỏ của cây đầy đủ.

\subsubsection{Tóm tắt và Nhận xét}
Mô hình Cây Quyết định với các tham số mặc định cho thấy hiệu suất dự đoán sản lượng nông nghiệp thấp hơn so với mô hình Hồi quy Tuyến tính trong nghiên cứu này, thể hiện qua chỉ số $R^2$ thấp hơn và các lỗi MAE, RMSE cao hơn. Đặc trưng quan trọng nhất được xác định là lượng phân bón sử dụng.
Dù vậy, việc phân tích độ quan trọng của đặc trưng và trực quan hóa cây cũng cung cấp những hiểu biết giá trị về dữ liệu.








\subsection{ Mô hình Mạng Nơ-ron Nhân tạo (ANN - MLPRegressor)}
\label{subsec:ann_results}

Để khám phá khả năng của các mô hình phức tạp hơn trong việc nắm bắt các mối quan hệ phi tuyến tiềm ẩn trong dữ liệu, nhóm em đã triển khai một Mạng Nơ-ron Nhân tạo dạng Multi-layer Perceptron Regressor (MLPRegressor). Một bước tiền xử lý quan trọng cho ANN là chuẩn hóa đặc trưng (feature scaling), do đó tất cả các đặc trưng đầu vào đã được chuẩn hóa bằng \texttt{StandardScaler} trước khi đưa vào mô hình. 

\subsubsection{Thông số và Kiến trúc Mô hình ANN}
\label{subsubsec:ann_architecture_metrics}

Mô hình MLPRegressor được cấu hình với các tham số chính như sau:
\begin{itemize}
    \item \textbf{Kiến trúc lớp ẩn (hidden\_layer\_sizes):} (100, 50) - bao gồm 2 lớp ẩn, lớp thứ nhất có 100 nơ-ron và lớp thứ hai có 50 nơ-ron.
    \item \textbf{Hàm kích hoạt (activation function):} ReLU (\texttt{relu}).
    \item \textbf{Thuật toán tối ưu (solver):} Adam (\texttt{adam}).
    \item \textbf{Số vòng lặp tối đa (max\_iter):} 500.
    \item \textbf{Dừng sớm (early\_stopping):} Được kích hoạt. Mô hình thực tế đã dừng sau 22 vòng lặp do không có sự cải thiện đáng kể trên tập validation.
\end{itemize}

Hiệu suất của mô hình ANN trên tập kiểm tra được đánh giá qua các chỉ số:
\begin{itemize}
    \item \textbf{Mean Absolute Error (MAE):} $0.5614$
    \item \textbf{Mean Squared Error (MSE):} $0.5047$
    \item \textbf{Root Mean Squared Error (RMSE):} $0.7104$
    \item \textbf{R-squared (R²):} $0.5218$
\end{itemize}

Giá trị $R^2 = 0.5218$ cho thấy mô hình ANN (với cấu hình ban đầu này) giải thích được khoảng $52.18\%$ sự biến thiên của sản lượng cây trồng. Kết quả này thấp hơn một chút so với mô hình Hồi quy Tuyến tính ($R^2 \approx 0.56$) nhưng cao hơn đáng kể so với mô hình Cây Quyết định với các tham số mặc định ($R^2 \approx 0.27$). Điều này cho thấy tiềm năng của ANN, tuy nhiên hiệu suất có thể được cải thiện thêm thông qua việc tinh chỉnh sâu hơn các siêu tham số và kiến trúc mạng.

Các chỉ số sai số MAE ($0.5614$) và RMSE ($0.7104$) cũng phản ánh mức độ lỗi tương tự như mô hình Hồi quy Tuyến tính.

\subsubsection{Trực quan hóa kết quả và quá trình huấn luyện}
\label{subsubsec:ann_visualization}

Ba biểu đồ chính được sử dụng để đánh giá mô hình ANN:

\begin{figure}[htb]
    \centering
    \includegraphics[width=0.8\textwidth]{Images/ann_actual_vs_predicted_plot.png}
    \caption{So sánh giá trị sản lượng thực tế và giá trị dự đoán bởi mô hình ANN trên tập kiểm tra.}
    \label{fig:ann_actual_vs_predicted}
\end{figure}
\FloatBarrier
Hình \ref{fig:ann_actual_vs_predicted} minh họa sự tương quan giữa giá trị sản lượng thực tế và giá trị dự đoán. Tương tự như mô hình Hồi quy Tuyến tính, các điểm dữ liệu có xu hướng tập trung quanh đường chéo $y=x$, cho thấy mô hình có khả năng dự đoán nhất định. Sự phân tán của các điểm phản ánh mức độ sai số của mô hình.

\begin{figure}[htb]
    \centering
    \includegraphics[width=0.8\textwidth]{Images/ann_residuals_distribution_plot.png} 
    \caption{Phân phối của phần dư (Actual - Predicted) từ mô hình ANN trên tập kiểm tra.}
    \label{fig:ann_residuals_distribution}
\end{figure}
\FloatBarrier
Phân phối phần dư của mô hình ANN (Hình \ref{fig:ann_residuals_distribution}) cũng tập trung quanh giá trị 0 và có hình dạng gần đối xứng, gợi ý rằng mô hình không mắc phải lỗi thiên lệch hệ thống nghiêm trọng.

\begin{figure}[htb]
    \centering
    \includegraphics[width=0.8\textwidth]{Images/ann_loss_curve.png}
    \caption{Đường cong hàm mất mát (Loss Curve) của mô hình ANN trong quá trình huấn luyện. Trục hoành thể hiện số vòng lặp (epochs).}
    \label{fig:ann_loss_curve}
\end{figure}
\FloatBarrier
Hình \ref{fig:ann_loss_curve} hiển thị giá trị của hàm mất mát (loss function) qua các vòng lặp huấn luyện. Có thể thấy giá trị hàm mất mát giảm nhanh trong những vòng lặp đầu tiên và sau đó hội tụ, cho thấy mô hình đã học được từ dữ liệu. Việc mô hình dừng sớm sau 22 vòng lặp (trong khi \texttt{max\_iter=500}) là do cơ chế \texttt{early\_stopping} được kích hoạt, giúp ngăn chặn overfitting và tiết kiệm thời gian huấn luyện khi không còn sự cải thiện đáng kể trên tập validation.

\subsubsection{Tóm tắt và Nhận xét}
Mô hình Mạng Nơ-ron Nhân tạo (MLPRegressor) với cấu hình ban đầu đã cho thấy hiệu suất dự đoán khá, với $R^2 \approx 0.52$, thấp hơn một chút so với Hồi quy Tuyến tính nhưng vượt trội hơn Cây Quyết định mặc định. Điều này cho thấy ANN có khả năng mô hình hóa các mối quan hệ trong dữ liệu, nhưng cũng nhấn mạnh tầm quan trọng của việc tinh chỉnh siêu tham số (như kiến trúc mạng, tốc độ học, phương pháp điều chuẩn) để khai thác tối đa tiềm năng của mô hình. Khác với các mô hình tuyến tính hay cây quyết định, việc diễn giải trực tiếp "độ quan trọng" của từng đặc trưng trong ANN phức tạp hơn và thường đòi hỏi các kỹ thuật phân tích bổ sung (ví dụ: Permutation Importance).








\subsection{Mô hình K-Nearest Neightbors (KNN Regressor)}
\label{subsec:knn_results}

Trong phần này, nhóm em trình bày kết quả của việc áp dụng mô hình K-Láng giềng Gần nhất (K-Nearest Neighbors - KNN) Regressor để dự đoán sản lượng cây trồng. Đây là một thuật toán dựa trên khoảng cách, do đó, tương tự như mô hình ANN, các đặc trưng đầu vào đã được nhóm em chuẩn hóa bằng \texttt{StandardScaler} trước khi huấn luyện để đảm bảo tính đồng nhất về thang đo.

\subsubsection{Thông số và hiệu suất mô hình KNN}
\label{subsubsec:knn_parameters_metrics}

Mô hình KNN Regressor được nhóm em cấu hình với các tham số chính sau:
\begin{itemize}
    \item \textbf{Số láng giềng gần nhất (K - n\_neighbors):} 5
    \item \textbf{Trọng số (weights):} `distance` (các láng giềng gần hơn có ảnh hưởng lớn hơn đến dự đoán)
    \item \textbf{Metric đo khoảng cách:} `minkowski` (với p=2, tương đương khoảng cách Euclidean)
\end{itemize}

Hiệu suất của mô hình KNN trên tập kiểm tra được nhóm em đánh giá qua các chỉ số:
\begin{itemize}
    \item \textbf{Mean Absolute Error (MAE):} $0.7771$
    \item \textbf{Mean Squared Error (MSE):} $0.9200$
    \item \textbf{Root Mean Squared Error (RMSE):} $0.9592$
    \item \textbf{R-squared (R²):} $0.1284$
\end{itemize}

Giá trị $R^2 = 0.1284$ cho thấy mô hình KNN (với $K=5$ và các thiết lập như trên) chỉ giải thích được khoảng $12.84\%$ sự biến thiên của sản lượng cây trồng. Đây là một kết quả tương đối thấp so với các mô hình khác mà nhóm em đã thử nghiệm trước đó (Hồi quy Tuyến tính $R^2 \approx 0.56$, ANN $R^2 \approx 0.52$). Điều này cho thấy mô hình KNN với cấu hình hiện tại chưa thực sự phù hợp hoặc chưa khai thác tốt các mối quan hệ trong bộ dữ liệu này.

Các chỉ số sai số MAE ($0.7771$) và RMSE ($0.9592$) cũng ở mức khá cao, phản ánh độ lệch lớn hơn giữa giá trị dự đoán và giá trị thực tế so với một số mô hình trước.

\subsubsection{Trực quan hóa kết quả mô hình KNN}
\label{subsubsec:knn_visualization}

Hai biểu đồ sau giúp nhóm em đánh giá trực quan hơn về hiệu suất của mô hình KNN.

\begin{figure}[htb]
    \centering
    \includegraphics[width=0.8\textwidth]{Images/knn_actual_vs_predicted_plot.png}
    \caption{So sánh giá trị sản lượng thực tế và giá trị dự đoán bởi mô hình KNN (K=5) trên tập kiểm tra. Hình ảnh do nhóm em tạo ra.}
    \label{fig:knn_actual_vs_predicted}
\end{figure}

Hình \ref{fig:knn_actual_vs_predicted} minh họa sự phân tán của các điểm dự đoán so với giá trị thực tế. Nhóm em nhận thấy rằng các điểm dữ liệu phân tán khá rộng và không tập trung chặt chẽ quanh đường chéo $y=x$. Điều này trực quan hóa cho kết quả $R^2$ thấp và các chỉ số lỗi cao đã được ghi nhận.

\begin{figure}[htb]
    \centering
    \includegraphics[width=0.8\textwidth]{Images/knn_residuals_distribution_plot.png}
    \caption{Phân phối của phần dư (Actual - Predicted) từ mô hình KNN (K=5) trên tập kiểm tra. Hình ảnh do nhóm em tạo ra.}
    \label{fig:knn_residuals_distribution}
\end{figure}

Phân phối phần dư của mô hình KNN (Hình \ref{fig:knn_residuals_distribution}) tuy có tâm lệch không quá xa 0, nhưng hình dạng của phân phối có vẻ rộng và không đều bằng một số mô hình khác, cho thấy sự biến thiên lớn trong sai số dự đoán.

\subsubsection{Tóm tắt và Nhận xét về Mô hình KNN}
Mô hình K-Láng giềng Gần nhất (KNN Regressor) với $K=5$ và các thiết lập ban đầu đã cho thấy hiệu suất dự đoán sản lượng nông nghiệp chưa cao trên bộ dữ liệu này. Kết quả $R^2$ thấp và các chỉ số lỗi MAE, RMSE cao hơn so với các mô hình như Hồi quy Tuyến tính và ANN.
Nhóm em cho rằng có một số nguyên nhân có thể dẫn đến kết quả này:
\begin{itemize}
    \item Giá trị $K=5$ có thể chưa phải là tối ưu. Việc lựa chọn $K$ phù hợp là rất quan trọng đối với KNN và thường cần được tinh chỉnh thông qua kiểm định chéo (Cross-Validation).
    \item KNN có thể nhạy cảm với "lời nguyền chiều dữ liệu" (curse of dimensionality) khi số lượng đặc trưng lớn (đặc biệt sau khi One-Hot Encoding).
    \item Bản chất của dữ liệu hoặc các mối quan hệ trong đó có thể không phù hợp với cách tiếp cận dựa trên láng giềng của KNN.
\end{itemize}
Các bước cải thiện tiềm năng cho mô hình KNN mà nhóm em có thể xem xét bao gồm việc tinh chỉnh siêu tham số $K$ một cách hệ thống, thử nghiệm các metric đo khoảng cách khác nhau, hoặc áp dụng các kỹ thuật giảm chiều dữ liệu trước khi đưa vào KNN. Tương tự như ANN, KNN cũng không cung cấp các hệ số hay độ quan trọng đặc trưng một cách trực tiếp như Hồi quy Tuyến tính hay Cây Quyết định.

% \newpage

% \section{So sánh và đánh giá hiệu suất các mô hình}
\label{sec:model_comparison_detailed}

Trong khuôn khổ của nghiên cứu này, bốn mô hình hồi quy đã được nhóm em xây dựng và đánh giá nhằm mục tiêu dự đoán sản lượng cây trồng (\texttt{Crop\_Yield\_MT\_per\_HA}) dựa trên các yếu tố về thời tiết, nông nghiệp, và các thông tin liên quan khác. Các mô hình bao gồm:
\begin{enumerate}
    \item Hồi quy Tuyến tính (Linear Regression)
    \item Cây Quyết định Regressor (Decision Tree Regressor) với các tham số mặc định
    \item Mạng Nơ-ron Nhân tạo (Artificial Neural Network - ANN) sử dụng MLPRegressor với cấu hình ban đầu
    \item K-nearest Neightbors (K-Nearest Neighbors - KNN) Regressor với K=5
\end{enumerate}
Tất cả các mô hình đều sử dụng cùng một quy trình tiền xử lý dữ liệu cơ bản (xử lý giá trị thiếu, mã hóa One-Hot cho biến phân loại). Đối với mô hình ANN và KNN, các đặc trưng đầu vào đã được chuẩn hóa bằng \texttt{StandardScaler} do tính nhạy cảm của các thuật toán này với thang đo của dữ liệu.

Phần này nhóm em sẽ trình bày so sánh chi tiết hiệu suất của các mô hình dựa trên các chỉ số đánh giá hồi quy phổ biến: R-squared (R²), Mean Absolute Error (MAE), Mean Squared Error (MSE), và Root Mean Squared Error (RMSE), được tính toán trên tập dữ liệu kiểm tra.

\subsection{Tổng hợp dết quả đánh giá hiệu suất}
\label{subsec:summary_evaluation_metrics_detailed}

Bảng \ref{tab:model_comparison_metrics_detailed} tóm tắt các chỉ số hiệu suất chính của từng mô hình. Giá trị R² phản ánh tỷ lệ phần trăm phương sai của biến mục tiêu mà mô hình có thể giải thích được. Các chỉ số MAE, MSE, và RMSE đo lường mức độ sai số trung bình của các dự đoán; giá trị càng thấp càng tốt.

\begin{table}[H] % Sử dụng [H] từ gói float để ép bảng ở đây (cẩn thận với bố cục)
    \centering
    \caption{Bảng so sánh chi tiết các chỉ số đánh giá hiệu suất của các mô hình hồi quy trên tập kiểm tra. Các giá trị lỗi (MAE, RMSE) được hiểu theo đơn vị của biến mục tiêu (ví dụ: tấn/ha).}
    \label{tab:model_comparison_metrics_detailed}
    \resizebox{\textwidth}{!}{% Bảng tự động co giãn theo chiều rộng trang
    \begin{tabular}{l c c c c}
        \toprule
        \textbf{Tên Mô hình} & \textbf{R-squared (R²)} & \textbf{MAE} & \textbf{MSE} & \textbf{RMSE} \\
        \midrule
        Hồi quy Tuyến tính & $0.5596$ & $0.5466$ & $0.4648$ & $0.6818$ \\
        Cây Quyết định Regressor (tham số gốc) & $0.2679$ & $0.6705$ & $0.7727$ & $0.8791$ \\
        Mạng Nơ-ron Nhân tạo (ANN - MLPRegressor)$^*$ & $0.5218$ & $0.5614$ & $0.5047$ & $0.7104$ \\
        K-nearest Neightbors (KNN, K=5)$^*$ & $0.1284$ & $0.7771$ & $0.9200$ & $0.9592$ \\
        \bottomrule
    \end{tabular}%
    }
    \footnotesize{$^*$ Các đặc trưng đầu vào cho mô hình ANN và KNN đã được chuẩn hóa bằng StandardScaler.}
\end{table}

\subsection{Phân tích chi tiết và thảo luận Kết quả}
\label{subsec:comparative_analysis_detailed}

Từ Bảng \ref{tab:model_comparison_metrics_detailed}, chúng tôi (trong báo cáo bạn sẽ dùng "nhóm em") tiến hành phân tích sâu hơn về hiệu suất của từng mô hình:

\subsubsection{Hồi quy Tuyến tính (Linear Regression)}
Mô hình Hồi quy Tuyến tính đạt được giá trị $R^2 = 0.5596$, cao nhất trong số bốn mô hình được thử nghiệm. Điều này có nghĩa là khoảng $55.96\%$ sự biến thiên trong sản lượng cây trồng có thể được giải thích bởi các đặc trưng đầu vào thông qua một mối quan hệ tuyến tính. Các chỉ số lỗi MAE ($0.5466$) và RMSE ($0.6818$) cũng là thấp nhất, cho thấy độ chính xác dự đoán tương đối tốt so với các mô hình khác trong thử nghiệm này.
\begin{itemize}
    \item \textbf{Ưu điểm:} Đơn giản, dễ triển khai, kết quả dễ diễn giải thông qua các hệ số hồi quy (như đã trình bày ở mục \ref{subsubsec:lr_coefficients}).
    \item \textbf{Hạn chế:} Giả định mối quan hệ tuyến tính giữa các đặc trưng và biến mục tiêu, có thể không nắm bắt được các tương tác phức tạp hoặc các mối quan hệ phi tuyến.
    \item \textbf{Đồ thị trực quan (Mục \ref{subsubsec:lr_visualization}):} Biểu đồ thực tế và dự đoán (Hình \ref{fig:lr_actual_vs_predicted}) cho thấy các điểm bám khá tốt quanh đường chéo, và phân phối phần dư (Hình \ref{fig:lr_residuals_distribution}) tương đối cân đối quanh 0.
\end{itemize}

\subsubsection{Mạng Nơ-ron nhân tạo (ANN - MLPRegressor)}
Mô hình ANN, với kiến trúc ban đầu gồm hai lớp ẩn (100 và 50 nơ-ron) và cơ chế dừng sớm (dừng sau 22 vòng lặp), đạt $R^2 = 0.5218$. Kết quả này rất đáng khích lệ, chỉ thấp hơn một chút so với Hồi quy Tuyến tính. MAE ($0.5614$) và RMSE ($0.7104$) cũng ở mức cạnh tranh.
\begin{itemize}
    \item \textbf{Ưu điểm:} Có khả năng học các mối quan hệ phi tuyến phức tạp mà không cần định nghĩa tường minh các đặc trưng tương tác hay đa thức. Đường cong hàm mất mát (Hình \ref{fig:ann_loss_curve}) cho thấy mô hình đã học hiệu quả.
    \item \textbf{Hạn chế:} Là một mô hình "hộp đen", khó diễn giải trực tiếp các trọng số. Hiệu suất rất nhạy cảm với kiến trúc mạng và các siêu tham số (tốc độ học, số epochs, hàm kích hoạt, thuật toán tối ưu, etc.), đòi hỏi quá trình tinh chỉnh kỹ lưỡng. Yêu cầu chuẩn hóa đặc trưng đầu vào.
    \item \textbf{Đồ thị trực quan (Mục \ref{subsubsec:ann_visualization}):} Các biểu đồ (Hình \ref{fig:ann_actual_vs_predicted} và \ref{fig:ann_residuals_distribution}) cũng cho thấy hiệu suất tốt, tương tự Hồi quy Tuyến tính.
\end{itemize}

\subsubsection{Cây quyết định Regressor (Decision Tree Regressor)}
Mô hình cây Quyết định với các tham số mặc định cho kết quả $R^2 = 0.2679$, thấp hơn đáng kể so với hai mô hình trên. Các chỉ số lỗi MAE ($0.6705$) và RMSE ($0.8791$) cũng cao hơn.
\begin{itemize}
    \item \textbf{Ưu điểm:} Dễ hiểu, dễ diễn giải thông qua cấu trúc cây (như Hình \ref{fig:dt_tree_visualization} đã minh họa một phần). Không yêu cầu chuẩn hóa đặc trưng. Có thể nắm bắt các mối quan hệ phi tuyến.
    \item \textbf{Hạn chế:} Một cây quyết định đơn lẻ (không giới hạn độ sâu hoặc không được cắt tỉa) rất dễ bị học quá khớp (overfitting) với dữ liệu huấn luyện, dẫn đến hiệu suất kém trên tập kiểm tra. Độ ổn định không cao, một thay đổi nhỏ trong dữ liệu có thể dẫn đến một cấu trúc cây rất khác.
    \item \textbf{Đồ thị trực quan (Mục \ref{subsubsec:dt_visualization}):} Biểu đồ thực tế và dự đoán (Hình \ref{fig:dt_actual_vs_predicted}) cho thấy sự phân tán lớn hơn và các "bậc thang" trong dự đoán, đặc trưng của cây quyết định.
\end{itemize}
Mặc dù độ quan trọng của đặc trưng (Mục \ref{subsubsec:dt_feature_importances}) cung cấp thông tin hữu ích, hiệu suất tổng thể cho thấy cần phải tinh chỉnh hoặc sử dụng các kỹ thuật ensemble dựa trên cây.

\subsubsection{K-nearest Neighbors (KNN Regressor)}
Mô hình KNN với $K=5$ và trọng số `distance` cho kết quả thấp nhất trong các thử nghiệm, với $R^2 = 0.1284$. MAE ($0.7771$) và RMSE ($0.9592$) cũng là cao nhất.
\begin{itemize}
    \item \textbf{Ưu điểm:} Đơn giản về mặt khái niệm, không có giả định mạnh về phân phối dữ liệu. Có thể hoạt động tốt với các mối quan hệ phi tuyến phức tạp cục bộ.
    \item \textbf{Hạn chế:} Hiệu suất rất phụ thuộc vào việc chọn giá trị $K$ và metric đo khoảng cách. Nhạy cảm với "lời nguyền chiều dữ liệu" (curse of dimensionality), đặc biệt khi có nhiều đặc trưng không liên quan hoặc dư thừa. Yêu cầu chuẩn hóa đặc trưng đầu vào. Chi phí tính toán dự đoán có thể cao với tập dữ liệu lớn do phải tính toán khoảng cách đến tất cả các điểm trong tập huấn luyện.
    \item \textbf{Đồ thị trực quan (Mục \ref{subsubsec:knn_visualization}):} Các biểu đồ (Hình \ref{fig:knn_actual_vs_predicted} và \ref{fig:knn_residuals_distribution}) phản ánh hiệu suất chưa cao của mô hình này.
\end{itemize}

\subsection{Thảo luận chung và Kết luận sơ bộ}
\label{subsec:overall_discussion_preliminary_conclusion}

Dựa trên các kết quả đánh giá ban đầu này, nhóm em kết luận:
\begin{enumerate}
    \item \textbf{Hồi quy Tuyến tính} đã thiết lập một baseline mạnh mẽ. Điều này cho thấy một phần đáng kể của mối quan hệ giữa các yếu tố đầu vào và sản lượng có thể là tuyến tính hoặc gần tuyến tính.
    \item \textbf{Mạng Nơ-ron Nhân tạo (ANN)} cho thấy tiềm năng lớn, đạt hiệu suất gần bằng Hồi quy Tuyến tính ngay cả với cấu hình ban đầu. Việc tinh chỉnh sâu hơn có thể giúp ANN vượt qua Hồi quy Tuyến tính, đặc biệt nếu có các mối quan hệ phi tuyến phức tạp mà mô hình tuyến tính bỏ lỡ.
    \item \textbf{Cây Quyết định và KNN}, với các thiết lập mặc định hoặc ban đầu, chưa cho thấy hiệu quả cao bằng hai mô hình trên. Cả hai mô hình này đều đòi hỏi sự tinh chỉnh siêu tham số cẩn thận (ví dụ: độ sâu của cây, số lượng mẫu tối thiểu để chia nhánh cho Cây Quyết định; giá trị $K$ cho KNN) để có thể đạt được hiệu suất tốt nhất.
\end{enumerate}

Kết quả này nhấn mạnh tầm quan trọng của việc lựa chọn mô hình phù hợp với bản chất của dữ liệu cũng như sự cần thiết của việc tinh chỉnh siêu tham số. Đối với các mô hình phức tạp như ANN, Cây quyết định, và KNN, việc sử dụng các giá trị tham số mặc định thường không mang lại kết quả tối ưu.



\end{document}


